\chapter{Worked Example: Lexical Ambiguity with Complex Types}

\newcommand{\ExNP}{\mathsf{np}}
\newcommand{\ExN}{\mathsf{n}}
\newcommand{\ExS}{\mathsf{s}}
\newcommand{\Parse}{P_{\Lex}}

\section{Purpose of the example}
The revised formalisation lets single words receive arbitrary Lambek formulae as lexical types.
This makes the examples much cleaner. Instead of introducing special atomic types for
intransitive verbs, transitive verbs, modifiers, or conjunctions, we use a small atomic inventory:
\[
\At_{\mathrm{ex}}=\{\ExNP,\ExN,\ExS\}.
\]
The familiar linguistic roles are then represented by complex formulae. For example,
\[
(\ExNP\backslash \ExS)/\ExNP
\]
is the type of a transitive verb, and
\[
(\ExS\backslash \ExS)/\ExS
\]
is a simple type for sentential conjunction.

\section{A toy corpus}
Let $\mathcal C_{\mathrm{ex}}$ consist of the following observed strings:
\[
\begin{array}{ll}
u_1 &= \text{cats nap},\\
u_2 &= \text{mice nap},\\
u_3 &= \text{cats chase mice},\\
u_4 &= \text{cats nap and mice nap},\\
u_5 &= \text{French wine drinker}.
\end{array}
\]
As before, $\Frag(\mathcal C_{\mathrm{ex}})$ contains the non-empty contiguous fragments of these
strings, such as \emph{cats chase}, \emph{chase mice}, \emph{and mice nap}, and
\emph{French wine drinker}.

\section{Lexical probabilities}
Assume the following probabilistic lexicon. Entries not listed have probability $0$.
\[
\begin{array}{c|c|c}
\text{word} & \text{formula} & \Lex(w,A)\\
\hline
\text{cats} & \ExNP & 0.95\\
\text{cats} & \ExN & 0.05\\
\text{mice} & \ExNP & 0.95\\
\text{mice} & \ExN & 0.05\\
\text{nap} & \ExNP\backslash \ExS & 1\\
\text{chase} & (\ExNP\backslash \ExS)/\ExNP & 1\\
\text{and} & (\ExS\backslash \ExS)/\ExS & 1\\
\text{French} & \ExN/\ExN & 0.45\\
\text{French} & (\ExN/\ExN)/(\ExN/\ExN) & 0.55\\
\text{wine} & \ExN/\ExN & 0.85\\
\text{wine} & \ExN & 0.15\\
\text{drinker} & \ExN & 1
\end{array}
\]
This table illustrates the main point of the revised formalisation. Words such as \emph{chase},
\emph{and}, \emph{French}, and \emph{wine} can be assigned non-atomic formulae directly. These
entries are lexical assignments, not values that must be obtained by recursively evaluating the
residual clauses.

\section{Simple sentence scores}
Consider first \emph{cats nap}. The lexical analysis
\[
\text{cats}:\ExNP,
\qquad
\text{nap}:\ExNP\backslash \ExS
\]
has lexical score
\[
0.95\cdot 1 = 0.95.
\]
It licenses a sentence because
\[
\ExNP\concat(\ExNP\backslash \ExS)\vdash_{\LC}\ExS.
\]
Hence
\[
\Parse(\text{cats nap},\ExS)\geq 0.95.
\]
In this toy lexicon there is no higher-scoring sentence analysis for the same string, so the parse
score is $0.95$.

For the transitive sentence \emph{cats chase mice}, the relevant lexical analysis is
\[
\text{cats}:\ExNP,
\qquad
\text{chase}:(\ExNP\backslash \ExS)/\ExNP,
\qquad
\text{mice}:\ExNP.
\]
The Lambek reduction is the usual one:
\[
\ExNP
\concat ((\ExNP\backslash \ExS)/\ExNP)
\concat \ExNP
\vdash_{\LC}
\ExNP\concat(\ExNP\backslash \ExS)
\vdash_{\LC}
\ExS.
\]
The score of this analysis is
\[
0.95\cdot1\cdot0.95=0.9025,
\]
so
\[
\Parse(\text{cats chase mice},\ExS)=0.9025.
\]

\section{Conjunction}
The word \emph{and} is assigned the complex formula
\[
(\ExS\backslash \ExS)/\ExS.
\]
For the string \emph{cats nap and mice nap}, the left substring \emph{cats nap} receives type
$\ExS$, and the right substring \emph{mice nap} also receives type $\ExS$. Then
\[
((\ExS\backslash \ExS)/\ExS)\concat \ExS
\vdash_{\LC}
\ExS\backslash \ExS,
\]
and therefore
\[
\ExS\concat(\ExS\backslash \ExS)\vdash_{\LC}\ExS.
\]
The full lexical score is
\[
0.95\cdot1\cdot1\cdot0.95\cdot1=0.9025,
\]
so the coordinated sentence receives
\[
\Parse(\text{cats nap and mice nap},\ExS)=0.9025.
\]

Notice what is \emph{not} happening here. We do not need to compute
\[
v(\text{and},(\ExS\backslash \ExS)/\ExS)
\]
from the recursive residual clause in order to say that \emph{and} has this type. The direct
lexical statement is
\[
\Lex(\text{and},(\ExS\backslash \ExS)/\ExS)=1.
\]
The recursive value $v(\text{and},(\ExS\backslash \ExS)/\ExS)$ would instead ask whether the
observed corpus contexts make \emph{and} behave like a right-looking residual expression. That
is a semantic-support question, not the lexical assignment itself.

\section{Ambiguity in \emph{French wine drinker}}
The phrase \emph{French wine drinker} now needs no special atomic types for origin, personhood,
or beveragehood. The ambiguity is represented entirely by different complex lexical entries.

The first reading is:
\begin{quote}
a French person who is a wine drinker.
\end{quote}
Here \emph{French} modifies the whole nominal \emph{wine drinker}. The lexical analysis is
\[
\text{French}:\ExN/\ExN,
\qquad
\text{wine}:\ExN/\ExN,
\qquad
\text{drinker}:\ExN.
\]
The reduction is
\[
(\ExN/\ExN)\concat((\ExN/\ExN)\concat \ExN)
\vdash_{\LC}
(\ExN/\ExN)\concat \ExN
\vdash_{\LC}
\ExN.
\]
The score is
\[
0.45\cdot0.85\cdot1=0.3825.
\]

The second reading is:
\begin{quote}
a drinker of French wine.
\end{quote}
Here \emph{French} modifies the nominal modifier \emph{wine}, producing the compound modifier
\emph{French wine}. The lexical analysis is
\[
\text{French}:(\ExN/\ExN)/(\ExN/\ExN),
\qquad
\text{wine}:\ExN/\ExN,
\qquad
\text{drinker}:\ExN.
\]
The reduction is
\[
((\ExN/\ExN)/(\ExN/\ExN))
\concat(\ExN/\ExN)
\concat \ExN
\vdash_{\LC}
(\ExN/\ExN)\concat \ExN
\vdash_{\LC}
\ExN.
\]
The score is
\[
0.55\cdot0.85\cdot1=0.4675.
\]

Thus both readings are grammatically available, but the second receives the higher score in this
toy lexicon:
\[
0.4675 > 0.3825.
\]
The ranking comes from the lexical probabilities, while the grammatical availability of each
reading comes from the Lambek reductions.

\section{Relation to the recursive valuation}
The examples above use $\Parse$, the lexical parse score. The recursive valuation $v$ still plays
the role defined in the previous chapter: it interprets formulae over observed corpus fragments
and supports the semantic consequence relation used in the soundness theorem.

The two layers answer different questions.
\begin{enumerate}[label=(\roman*)]
    \item $\Lex(w,A)$ asks whether a word is assigned type $A$ in the lexicon.
    \item $\Parse(x,A)$ asks how strongly a string can be parsed as type $A$ using lexical entries and
    Lambek derivability.
    \item $v(x,A)$ asks how strongly an observed fragment behaves as an instance of $A$ under the
    recursive product and residual clauses.
\end{enumerate}
This separation is useful because it lets the lexicon contain ordinary complex Lambek types while
leaving the recursive semantic clauses clean enough for the soundness proof.
